% !TeX root = ../../MainBook.tex
\section{信号与系统简介}\label{sec:Introduction}

\subsection{导言}\label{sub:basic_concepts}

《信号与系统》研究的是确定信号通过线性时不变系统的过程。信号是信息的表现形式与传送载体，信息是蕴含在信号中的具体内容。如何处理信号、如何分析系统这两个问题形成了本书的主线。

本小节中，我们首先对各章的内容进行简要的总结，以供读者根据需要进行阅读。

第一章给出了信号与系统课程中的一些基本定义，并补充了一些本书要用到的数学基础，对工科的教学内容进行了一些扩展。其中标“*”号的小节，主要是为了逻辑的完整性，多数读者可以跳过它们。实变函数实际上是传统教材中“信号空间”所介绍内容的背景，读者应重点关注我们介绍的几个线性赋范空间。我们将看到，用实变函数的语言来描述信号是准确而高效的。

第二章介绍了傅里叶变换的基本内容。傅里叶变换是信号分析的基石，从变换域的角度讨论周期函数、非周期函数的合成与分解，能够带来新的认识，是单纯的时域分析无法做到的。我们在这一章给出简要的分布理论，它能够解释课程中涉及到的一些形式计算的原理。

第三章给出了连续系统的概述，介绍了什么是系统的线性、时不变性，以及一系列其他的系统特性。我们所关注的系统，往往是由特殊类型的微分方程描述的，从系统的视角出发，可以得到一些经典视角下难以发觉的分析方法。其中，微分方程组实际上涉及到多入多出系统，初学者不妨略过这一节。

第四章中，我们结合专业需求对傅里叶变换进行了一些更深入的讨论。采样与插值问题是信号处理的核心问题之一，采用傅里叶变换的分析方法，可以得到对理想采样模型最深入的认识。多维傅里叶变换在图像信号处理以及半导体器件中都有着重要的应用。我们还将通过傅里叶变换分析其他变换。此外，对于众多文献中不加说明的“有界区间傅里叶变换”和“有界区间卷积”，本章的最后一节讨论了其原理，还建立了初步的小波分析法。

第五章介绍了拉普拉斯变换，它可以看作是傅里叶变换的推广。通过拉普拉斯变换，我们能够加深对系统的认识。对于多数教材没有详细讲解的留数定理方法与确定收敛域的方法，这一章中进行了详细的讨论。

第六章中我们开始讨论离散信号。z变换是拉普拉斯变换的离散版，而离散时间傅里叶变换是傅里叶变换的离散版。这一章最后，我们将讨论离散傅里叶变换并指出它如何帮助人们用计算机实现信号的变换域分析。

第七章是对离散系统的讨论。这一章的多数内容基本就是第三章内容的离散版本。其中的差分方程组涉及到多入多出系统，本书暂未介绍，但不妨碍对其进行分析。尽管连续系统涉及更加复杂的数学背景，但考虑到工科生的组合数学基础相对薄弱，而一些边界条件的处理又较为繁琐，本章单列于本书的后半部分。读者可根据自身需求安排阅读顺序。

在附录中，我们又一次讨论了傅里叶变换，这一章中的内容更加严谨与困难，因此没有直接并入第二章中，而是作为补充内容放在本书的最后。它补全了第二章的逻辑。这一章仅供感兴趣的读者参考。

\subsection{连续信号与离散信号}\label{sub:signal}
信号的本质是函数，\textbf{连续信号}就是定义域“连续”的函数，例如一个时间的函数$f(t)$，一般要求定义域具有连续统的势。相应的，\textbf{离散信号}的定义域往往是至多可数集，于是我们常认为其定义域是整数集，并记离散信号为$x[n]$，用研究数列的方法研究离散信号。

离散信号可以视作连续信号的定义域离散化的结果，一般来说，这样一个过程通过对连续信号进行采样来实现；而数字信号则是离散信号的幅值离散化的结果，是对离散信号进行量化的结果。量化的概念以及数字信号将在数字信号处理的课程中着重讨论。

鉴于信号与函数、离散信号与序列在数学本质上一致，本书在数学分析与变换理论的讨论中将分别统一称为“函数”和“序列”，而在涉及系统分析与工程背景的专有术语中仍保留“信号”的称呼。

从电学的经验来看，电压、电流信号的功率总为其\textbf{幅值的平方}乘以某个常数，于是对于一般的信号我们也采取此定义：

\begin{definition}[能量]
    连续信号$f:\mathbb{R}\to\mathbb{C}$的\textbf{能量}(energy)定义为
    \begin{equation}
        E_f=\int_{-\infty}^{\infty}|f(x)|^2\,\dd x,
    \end{equation}
    对于序列$x:\mathbb{Z}\to\mathbb{C}$，
    \begin{equation}
        E_x=\sum_{n = -\infty}^{\infty} |x[n]|^2.
    \end{equation}
\end{definition}

真实的物理世界中的信号的能量总是有限的，但是，为了更好地研究它们，很多时候使用真实信号的一部分做适当延拓作为研究对象，因此我们将遇到一些不满足“平方可积”性质$\int_{-\infty}^{\infty}|f(x)|^2\,\dd x<\infty$的函数。例如，幅值不为零的正弦函数在整个时间轴上的能量为无穷，但在足够长的时间内，能量与观察时长近似成正比。为了描述这类信号，需要使用平均功率。记长度为$T$的对称窗口上的平均功率为
\begin{equation}
    P_T(f)=\frac{1}{T}\int_{-T/2}^{T/2}|f(t)|^2\,\dd t.
\end{equation}
若$T\to\infty$的极限存在，则定义
\begin{equation}
    P_f=\lim_{T\to\infty}P_T(f).
\end{equation}
对于序列也有类似的定义：
\begin{equation}
    P_x=\lim_{M\to\infty}\frac{1}{2M+1}\sum_{n=-M}^M |x[n]|^2.
\end{equation}
\begin{definition}[能量信号与功率信号]
    能量非零有限的信号称为\textbf{能量信号}：
    \begin{equation}
        0<\int_{-\infty}^{\infty}|f(x)|^2\,\dd x<\infty,
    \end{equation}
    能量无限而功率非零有限的信号称为\textbf{功率信号}：
    \begin{equation}
        \int_{-\infty}^{\infty}|f(x)|^2\,\dd x=\infty,\quad \lim_{T \to \infty}\frac{1}{T}\int_{-T/2}^{T/2}|f(x)|^2\,\dd x<\infty.
    \end{equation}
\end{definition}

能量无限等价于功率非零。在\refsub{Lp_space}中我们将看到，能量信号是$L^2$空间中的元素。我们还可以将这一定义推广到序列，其中离散的能量信号满足“平方可和”性质$\sum_{n=-\infty}^{\infty}|x[n]|^2<\infty$，对应$l^2$空间中的元素，并将能量无限而功率有限的序列称为功率信号。

\subsection{典型的信号}\label{sub:ordinary_signal}
\subsubsection{连续信号}
\begin{definition}[采样函数]\label{def:sinc}
    \begin{equation}
        \Sa(t)=\frac{\sin(t)}{t}.
    \end{equation}
\end{definition}
它在0处连续延拓为1，在$\pi$的整数倍处为0，是偶函数，在正半轴上的积分值为$\pi /2$（这个结果称为Dirichlet积分）。国外教材中采样函数一般为一个类似的函数
\begin{equation}
    \sinc(t)=\frac{\sin(\pi t)}{\pi t}=\Sa(\pi t).
\end{equation}
它们的参数不同但性质类似、本质相同。需要注意的是，在绘图时python和MATLAB中只有$\sinc$函数，而且实际上指的是$\Sa$函数。
\begin{definition}[钟形函数/高斯函数]
    \begin{equation}
        f(t)=E\e^{-(t/\tau)^2},
    \end{equation}
    $\tau$为衰减速率或时间常数，$\tau$越大，函数值衰减越慢。
\end{definition}
\begin{definition}[单位脉冲函数/Dirac函数]\label{def:delta}
    \begin{equation}
        \delta(x)=\begin{dcases}
            0,      & x\neq 0, \\
            \infty, & x=0
        \end{dcases}
    \end{equation}
    满足：\\
    \ding{172} 采样性质：$f(x)\delta(x)=f(0)\delta(x)$；\\
    \ding{173} 积分性质：$\int_{-\infty}^{\infty}\delta(x)\,\dd x=1$，从而$\int_{-\infty}^{\infty}f(x)\delta(x)\,\dd x=f(0)$；\\
    \ding{174} $\delta(x)$函数是偶函数：$\delta(-x)=\delta(x)$。
\end{definition}
这个函数可视为一些性质较好的函数如高斯函数、采样函数集中到$x=0$附近时的极限情况，从而具备很多优美的性质，见\refsub{approach}。实际上，它应该作为一个分布（而不是函数）来理解，并且是无限阶可导的，我们将在\refsub{general_distributions}中详细讨论这一点。目前，我们仅能从形式和直观上定义这样一个函数。

根据以上定义中列举的性质，可以得到$\delta$函数的伸缩变换性质：$\displaystyle \delta(ax)=\frac{1}{|a|}\delta(x)$，这是因为我们需要保证：
\begin{equation}
    \int_{-\infty}^{\infty}\delta(ax)\,\dd(ax)=a\int_{-\infty}^{\infty}\delta(y)\,\dd y=a.
\end{equation}
我们规定，$\delta_a(x)=\delta(x-a)$，即在$x=a$处采样的Dirac函数。

\begin{definition}[单位阶跃函数]
    \begin{equation}
        u(x)=\begin{dcases}
            0, & x<0, \\
            1, & x>0.
        \end{dcases}
    \end{equation}
\end{definition}

通过对$x$的正负进行分类讨论，可以得到含单位阶跃函数的积分公式，后面会多次使用它：
\begin{equation}
    \int_{-\infty}^{x}u(y)f(y)\,\dd y=u(x)\int_{0}^{x}f(y)\,\dd y.
\end{equation}

\begin{figure}[H]
    \centering
    \includegraphics[width=0.8\textwidth]{Figure_1}
    \caption{采样函数、高斯函数、Dirac函数、单位阶跃函数}
\end{figure}

这个函数在0处的值可以任意定义，一般取0、1或$1/2$，不会有影响，所以一般也不会讨论。

\begin{definition}
    设集合$A\subset\mathbb{R}$，其\textbf{特征函数}(indicator function)为
    \begin{equation}
        \chi_A(x)=\begin{dcases}
            1, & x\in A,    \\
            0, & x\notin A.
        \end{dcases}
    \end{equation}
\end{definition}

例如，一个幅值为1、存在于区间$[a,b]$上的矩形波可以表示为$\chi_{[a,b]}(x)$。在多数信号与系统的教材中，这类函数通常会使用其他的符号来定义，例如$R_T(t),\ G_T(t)$等，但都不如特征函数的符号意义清晰。

我们将在实变函数的理论中用到它。例如，Dirichlet函数$D(x)$可表示为$\chi_{\mathbb{Q}}(x)$。概率论中会定义连续性随机变量的特征函数(characteristic function)，指随机变量的矩母函数的傅里叶变换，本书中不会用到它，请读者注意它们的区别。

单位阶跃函数的积分为单位斜变函数(the unit ramp)，导数为Dirac函数，从直观上这不难理解，在学习\refsub{general_distributions}之后，很容易自行验证它。

\begin{definition}[单边指数函数和双边指数函数]\label{def:exp}
    \begin{equation}
        f(t)=\e^{-at}u(t)=\begin{dcases}
            \e^{-at}, & \text{if }t>0, \\
            0,        & \text{if }t<0
        \end{dcases}
    \end{equation}
    称为单边指数函数，
    \begin{equation}
        g(t)=f(t)+f(-t)=\begin{dcases}
            \e^{-at}, & \text{if }t\geq 0, \\
            \e^{at},  & \text{if }t<0
        \end{dcases}
    \end{equation}
    称为双边指数函数。
\end{definition}

\subsubsection{序列}
\begin{definition}[单位脉冲序列/Kronecker $\delta$函数]
    \begin{equation}
        \delta[n]=\begin{dcases}
            1, & \text{if }n=0,     \\
            0, & \text{if }n\neq 0.
        \end{dcases}
    \end{equation}
\end{definition}
将它进行时移，可得
\begin{equation}
    \delta_k[n]=\delta[n-k]=\begin{dcases}
        1, & \text{if }n=0,     \\
        0, & \text{if }n\neq 0.
    \end{dcases}
\end{equation}
它与Dirac $\delta$函数一样，都具有采样性质：$x[n]\delta_k[n]=x[k]\delta_k[n],\ \forall n$。
\begin{definition}[单位阶跃序列]
    \begin{equation}
        u[n]=\begin{dcases}
            1, & \text{if }n\geq 0, \\
            0, & \text{if }n< 0.
        \end{dcases}
    \end{equation}
\end{definition}
类似连续信号，由它可以衍生出矩形窗序列：
\begin{equation}
    R_N[n]=u[n]-u[n-N]=\begin{dcases}
        1, & \text{if }0\leq n\leq N-1, \\
        0, & \otherwise.
    \end{dcases}
\end{equation}

\begin{definition}[指数序列]
    $x[n]=a^n u[n]$称为指数序列，其中$a\in\mathbb{R}$。
\end{definition}
\begin{definition}[复指数序列]
    $E_{\Omega}[n]=\e^{\im\Omega_0 n}$称为复指数序列，其中$\Omega_0$称为数字角频率。
\end{definition}

$E_{\Omega}[n]$这个记号仅在本书中使用，但这样做是值得的——它表征了序列频率成分的基本形式，我们将在序列的章节中多次用到它。

\begin{definition}[正弦序列]
    $x[n]=\sin(\Omega_0 n)$，其中$\Omega_0$称为数字角频率。
\end{definition}
\begin{figure}[H]
    \centering
    \includegraphics[width=0.5\textwidth]{disct_sin}
    \caption{正弦序列}
\end{figure}

\begin{definition}[符号序列]
    \begin{equation}
        x[n]=\sgn[n]=\begin{dcases}
            -1, & \text{if }n<0, \\
            0,  & \text{if }n=0, \\
            1,  & \text{if }n>0.
        \end{dcases}
    \end{equation}
\end{definition}

下面介绍序列的伸缩变换：
\begin{equation}
    x[n]\to x[\alpha n],\quad \alpha\in\mathbb{R}.
\end{equation}
如果$\alpha\in\mathbb{N}_+$，则伸缩后的序列仍为序列，只是遗漏了原序列的一些信息；如果$\alpha=\frac{1}{N},\ N\in\mathbb{N}_+$，则伸缩后的序列在$N$的倍数以外的点没有定义，我们必须人为规定它们的值，这就是\textbf{插值}。

最简单的处理方式是将这些点全部赋值为0。我们定义：
\begin{definition}[序列的伸缩]
    序列的伸缩分为两种：
    \begin{equation}
        y[n]=x[Mn],\quad M\in\mathbb{N}_+
    \end{equation}
    称为\textbf{降采样/下采样}，对应\textbf{抽取}；
    \begin{equation}
        y[n]=\begin{dcases}
            x\sbr{\frac{n}{N}}, & n=Nk,\ N\in\mathbb{Z}, \\
            0,                  & \otherwise
        \end{dcases}
    \end{equation}
    称为\textbf{升采样/上采样}，对应\textbf{插值}。
\end{definition}

如果序列本身是由连续信号采样得到的，例如$x[n]=f(T_s n)$，那么插值就对应于还原连续信号的过程，因此序列的伸缩相当于处理连续信号时增加或减少样值点。降采样、升采样的名称就来源于此。对于降采样、升采样，我们将在\refsub{interpolation}中详细讨论。
